% 03-monitor-control.tex - Monitor and Control Interface

\chapter{Monitor and Control Interface}
\label{sec:monitor-control}

This section defines the command interface for \SC{}, including the UDP packet protocol, system commands, data management commands, and MIB entries.

\begin{noteBox}
\SC{} uses a UDP packet format based on the \MCS{} standard and shares the same network infrastructure as other LWA subsystems, but operates independently. Like HAL, \MCS{} does not see or route \SC{} messages, and \SC{} cannot see \MCS{} or subsystem messages. However, unlike HAL, \SC{} packets are not cryptographically signed.
\end{noteBox}

\section{UDP Command Protocol}

\SC{} uses UDP packets for its command interface, following a packet format based on the \MCS{} standard.

\subsection{Network Configuration}

The default network port assignments are listed in Table~\ref{tab:network-ports}.

\begin{table}[htbp]
    \centering
    \caption{Default Network Ports}
    \label{tab:network-ports}
    \begin{tabular}{lcl}
        \toprule
        \textbf{Port} & \textbf{Default} & \textbf{Description} \\
        \midrule
        Receive Port & 5050 & Incoming commands to \SC{} \\
        Send Port & 5051 & Outgoing responses from \SC{} \\
        Reference Port & 5052 & \O{}MQ REP socket for reference IDs \\
        \bottomrule
    \end{tabular}
\end{table}

Port assignments are configurable via the \file{defaults.json} configuration file using the \code{mcs.message\_in\_port}, \code{mcs.message\_out\_port}, and \code{mcs.message\_ref\_port} parameters.

The server bind address is automatically determined by scanning network interfaces for a \code{10.1.x.0} network. The bind address is set to \code{10.1.x.2} on the matching subnet.

\subsection{Packet Structure}

Command packets use the following structure, based on the \MCS{} packet format (Table~\ref{tab:packet-structure}).

\begin{table}[htbp]
    \centering
    \caption{Command Packet Structure}
    \label{tab:packet-structure}
    \begin{tabular}{cclp{5.5cm}}
        \toprule
        \textbf{Bytes} & \textbf{Length} & \textbf{Field} & \textbf{Description} \\
        \midrule
        0--2 & 3 & Destination & Subsystem identifier (``SCM'') \\
        3--5 & 3 & Sender & Source subsystem identifier \\
        6--8 & 3 & Command & Three-letter command code \\
        9--17 & 9 & Reference & Reference number (zero-padded integer) \\
        18--21 & 4 & Data Length & Length of data section \\
        22--27 & 6 & MJD & Modified Julian Date \\
        28--36 & 9 & MPM & Milliseconds Past Midnight \\
        37 & 1 & - & Space \\
        38+ & Variable & Data & Command-specific data payload \\
        \bottomrule
    \end{tabular}
\end{table}

\subsection{Response Format}

Response packets follow a similar structure (Table~\ref{tab:response-structure}).

\begin{table}[htbp]
    \centering
    \caption{Response Packet Structure}
    \label{tab:response-structure}
    \begin{tabular}{cclp{5.5cm}}
        \toprule
        \textbf{Bytes} & \textbf{Length} & \textbf{Field} & \textbf{Description} \\
        \midrule
        0--2 & 3 & Destination & Original sender \\
        3--5 & 3 & Sender & ``SCM'' \\
        6--8 & 3 & Command & Echoed from request \\
        9--17 & 9 & Reference & Echoed from request \\
        18--21 & 4 & Data Length & Response data length + 8 \\
        22--27 & 6 & MJD & Current timestamp \\
        28--36 & 9 & MPM & Current timestamp \\
        37 & 1 & - & Space \\
        38 & 1 & Response Code & \code{A} = Accepted, \code{R} = Rejected \\
        39--45 & 7 & Status & System status string \\
        46+ & Variable & Data & Response data \\
        \bottomrule
    \end{tabular}
\end{table}

\subsection{Reference ID Server}

\SC{} provides a \O{}MQ REP socket on the reference port (default: 5052) for generating unique reference IDs. Clients send the message \code{next\_ref} and receive a unique integer in response. Reference IDs:

\begin{itemize}
    \item Range from 1 to 999,999,999
    \item Roll over to 1 when the maximum is reached
    \item Are persisted to the file \file{.sc\_reference\_id} every 10 IDs
    \item Resume from the persisted value (plus 10) after a restart
\end{itemize}

\section{System Commands}

The system commands are listed in Table~\ref{tab:system-commands}.

\begin{table}[htbp]
    \centering
    \caption{System Commands}
    \label{tab:system-commands}
    \begin{tabular}{lp{6cm}p{4cm}}
        \toprule
        \textbf{Command} & \textbf{Description} & \textbf{Parameters} \\
        \midrule
        \cmd{PNG} & Ping/heartbeat verification & None \\
        \cmd{INI} & Initialize \SC{} system & None \\
        \cmd{SHT} & Shutdown \SC{} system & None \\
        \cmd{RPT} & Report MIB entry value & MIB key string \\
        \bottomrule
    \end{tabular}
\end{table}

\subsection{PNG --- Ping}

The \cmd{PNG} command verifies system responsiveness. No parameters are required. Returns an accepted response if the system is operational.

\subsection{INI --- Initialize}

The \cmd{INI} command initializes the \SC{} system:
\begin{itemize}
    \item Transitions system state to \status{BOOTING}
    \item Creates or restarts the station monitoring thread
    \item Creates or restarts the per-\DR{} queue manager threads (DR1--DR5)
    \item Restores pending queue items from the SQLite database
    \item Clears the global inhibit flag for all \DR{} copy queues
    \item Transitions system state to \status{NORMAL} upon completion
\end{itemize}

The initialization process runs asynchronously. The command is rejected with exit code \hex{03} if another \cmd{INI} or \cmd{SHT} operation is already in progress.

\subsection{SHT --- Shutdown}

The \cmd{SHT} command shuts down the \SC{} system:
\begin{itemize}
    \item Transitions system state to \status{SHUTDWN} immediately
    \item Pauses all \DR{} copy queues
    \item Stops all \DR{} queue manager threads
    \item Stops the station monitoring thread
\end{itemize}

The shutdown process runs asynchronously. The command is rejected with exit code \hex{03} if another \cmd{INI} or \cmd{SHT} operation is already in progress. Pending queue items are preserved in the SQLite database and will be restored on the next \cmd{INI}.

\subsection{RPT --- Report}

The \cmd{RPT} command queries MIB entries. The data field contains the MIB key string to query. See Section~\ref{sec:mib-entries} for the full list of queryable entries.

\section{Data Management Commands}

The data management commands are listed in Table~\ref{tab:data-commands}.

\begin{table}[htbp]
    \centering
    \caption{Data Management Commands}
    \label{tab:data-commands}
    \begin{tabular}{lp{5cm}p{5cm}}
        \toprule
        \textbf{Command} & \textbf{Description} & \textbf{Parameters} \\
        \midrule
        \cmd{SCP} & Queue a file copy & \code{host:path->dest:path} \\
        \cmd{SRM} & Queue a file delete & \code{[-tNOW] host:path} \\
        \cmd{PAU} & Pause a \DR{} queue & \code{DR\{1-5\}} or \code{ALL} \\
        \cmd{RES} & Resume a \DR{} queue & \code{DR\{1-5\}} or \code{ALL} \\
        \cmd{SCN} & Cancel a queued copy & Copy ID string \\
        \bottomrule
    \end{tabular}
\end{table}

\subsection{SCP --- Smart Copy}

The \cmd{SCP} command queues a file copy operation. The data field contains the source and destination in the format:

\begin{lstlisting}
host:path->dest:path
\end{lstlisting}

Where \param{host} is the source hostname (or \DR{} name), \param{path} is the full source file path, \param{dest} is the destination hostname, and the second \param{path} is the destination path. Colons in file paths must be escaped with a backslash.

On success, the response data contains the assigned copy ID. The \DR{} is determined from the \param{host} field. The command is rejected with:
\begin{itemize}
    \item \hex{04} if the system is not initialized
    \item \hex{01} if the \DR{} name is not recognized
    \item \hex{02} if an error occurs while queuing the copy
\end{itemize}

\subsection{SRM --- Smart Remove}

The \cmd{SRM} command queues a file delete operation. The data field format is:

\begin{lstlisting}
[-tNOW] host:path
\end{lstlisting}

Where \param{host} is the hostname and \param{path} is the file to delete. The optional \code{-tNOW} flag causes the deletion to be executed immediately rather than deferred to the next purge cycle. Without \code{-tNOW}, the file is marked for deletion but only removed during the daily purge.

On success, the response data contains the assigned delete ID. The command is rejected with:
\begin{itemize}
    \item \hex{04} if the system is not initialized
    \item \hex{01} if the \DR{} name is not recognized
    \item \hex{02} if an error occurs while queuing the delete
\end{itemize}

\subsection{PAU --- Pause}

The \cmd{PAU} command pauses copy queue processing for a specified \DR{}. The data field contains either a specific \DR{} name (e.g., \code{DR1}) or \code{ALL} to pause all queues. When paused:

\begin{itemize}
    \item The currently active copy is interrupted (the \code{rsync} process is killed)
    \item No new items are pulled from the queue
    \item The global inhibit flag is set, preventing the station monitor from auto-resuming the queue
\end{itemize}

The command is rejected with:
\begin{itemize}
    \item \hex{04} if the system is not initialized
    \item \hex{01} if the \DR{} name is not recognized
    \item \hex{02} if an error occurs while pausing the queue
\end{itemize}

\subsection{RES --- Resume}

The \cmd{RES} command resumes copy queue processing for a specified \DR{}. The data field contains either a specific \DR{} name or \code{ALL}. When resumed:

\begin{itemize}
    \item The currently paused copy is restarted (using \code{rsync --append} for seamless continuation)
    \item Queue processing resumes from where it was paused
    \item The global inhibit flag is cleared, allowing the station monitor to manage the queue
\end{itemize}

The command is rejected with:
\begin{itemize}
    \item \hex{04} if the system is not initialized
    \item \hex{01} if the \DR{} name is not recognized
    \item \hex{02} if an error occurs while resuming the queue
\end{itemize}

\subsection{SCN --- Smart Cancel}

The \cmd{SCN} command cancels a previously queued copy or delete operation. The data field contains the copy ID string returned by the \cmd{SCP} or \cmd{SRM} command. If the specified copy is currently active, it is stopped. The item is marked as canceled in the results cache and will be skipped when encountered in the queue.

The command is rejected with:
\begin{itemize}
    \item \hex{04} if the system is not initialized
    \item \hex{01} if the copy ID is not found
    \item \hex{02} if an internal error occurs
\end{itemize}

\section{MIB Entries}
\label{sec:mib-entries}

The MIB entries listed in Table~\ref{tab:mib-entries} are queryable via the \cmd{RPT} command. All MIB values are returned as ASCII text strings.

\begin{table}[htbp]
    \centering
    \caption{General MIB Entries}
    \label{tab:mib-entries}
    \begin{tabular}{lp{8cm}}
        \toprule
        \textbf{Entry} & \textbf{Description} \\
        \midrule
        \mibentry{SUMMARY} & Current system status (7-character string) \\
        \mibentry{INFO} & Informational message (up to 256 characters) \\
        \mibentry{LASTLOG} & Last log entry (up to 256 characters) \\
        \mibentry{SUBSYSTEM} & Subsystem identifier (``SCM'') \\
        \mibentry{SERIALNO} & Serial number \\
        \mibentry{VERSION} & Software version string \\
        \bottomrule
    \end{tabular}
\end{table}

\subsection{Per-DR Observing Status}

The observing status of each \DR{} is available via entries of the form \mibentry{OBSSTATUS\_\{DR\}} (Table~\ref{tab:obsstatus-mib}).

\begin{table}[htbp]
    \centering
    \caption{Observing Status MIB Entry}
    \label{tab:obsstatus-mib}
    \begin{tabular}{lp{8cm}}
        \toprule
        \textbf{Entry} & \textbf{Description} \\
        \midrule
        \mibentry{OBSSTATUS\_DR\{n\}} & Recording state of \DR{}\{n\} (``busy'' or ``idle'') \\
        \bottomrule
    \end{tabular}
\end{table}

Example: \mibentry{OBSSTATUS\_DR3} returns the recording state of DR3.

\subsection{Per-DR Queue Status}

Queue information for each \DR{} is available via entries of the form \mibentry{QUEUE\_\{property\}\_\{DR\}} (Table~\ref{tab:queue-mib}).

\begin{table}[htbp]
    \centering
    \caption{Queue MIB Entries}
    \label{tab:queue-mib}
    \begin{tabular}{lp{8cm}}
        \toprule
        \textbf{Entry} & \textbf{Description} \\
        \midrule
        \mibentry{QUEUE\_SIZE\_DR\{n\}} & Number of pending items in the queue (including active copy) \\
        \mibentry{QUEUE\_STATUS\_DR\{n\}} & Queue state (``active'', ``paused'', or ``global paused'') \\
        \mibentry{QUEUE\_STATS\_DR\{n\}} & Queue statistics: pending, processing, completed, and failed counts \\
        \mibentry{QUEUE\_ENTRY\_\{id\}} & Status of a specific copy command by ID \\
        \bottomrule
    \end{tabular}
\end{table}

Example: \mibentry{QUEUE\_SIZE\_DR1} returns the number of items queued for DR1.

\subsection{Per-DR Active Copy Status}

Information about the currently active copy on each \DR{} is available via entries of the form \mibentry{ACTIVE\_\{property\}\_\{DR\}} (Table~\ref{tab:active-mib}).

\begin{table}[htbp]
    \centering
    \caption{Active Copy MIB Entries}
    \label{tab:active-mib}
    \begin{tabular}{lp{8cm}}
        \toprule
        \textbf{Entry} & \textbf{Description} \\
        \midrule
        \mibentry{ACTIVE\_ID\_DR\{n\}} & Copy ID of the active transfer (``None'' if idle) \\
        \mibentry{ACTIVE\_STATUS\_DR\{n\}} & Status string of the active transfer \\
        \mibentry{ACTIVE\_BYTES\_DR\{n\}} & Bytes transferred by the active copy \\
        \mibentry{ACTIVE\_PROGRESS\_DR\{n\}} & Percentage progress (e.g., ``42\%'') \\
        \mibentry{ACTIVE\_SPEED\_DR\{n\}} & Current transfer speed (e.g., ``12.34MB/s'') or ``paused'' \\
        \mibentry{ACTIVE\_REMAINING\_DR\{n\}} & Estimated time remaining (e.g., ``01:23:45'') or ``unknown'' \\
        \bottomrule
    \end{tabular}
\end{table}

Example: \mibentry{ACTIVE\_PROGRESS\_DR2} returns the percentage progress of the active copy on DR2.
